\documentclass[12pt, a4paper]{article}

% -----------------------------------------------------------
% Packages
% -----------------------------------------------------------

\usepackage[utf8]{inputenc} % UTF-8 encoding
\usepackage[margin=1in]{geometry} % Page geometry
\usepackage{amsmath, amssymb} % Math symbols and equations
\usepackage{graphicx} % Include images
\usepackage{hyperref} % Clickable links in PDF
\usepackage{setspace} % Line spacing
\usepackage{xcolor} % Text color
\usepackage{titlesec} % Custom section titles
\usepackage{enumitem} % Custom lists

% -----------------------------------------------------------
% Document Configuration
% -----------------------------------------------------------

\title{Fisica - domande teoria}
\author{Claudio Tessa}
\date{\the\year}
    
% Custom section styling
\titleformat{\section}
  {\normalfont\large\bfseries}{\thesection}{1em}{}

% -----------------------------------------------------------
% Begin Document
% -----------------------------------------------------------

\begin{document}
\maketitle


\newcounter{qn} % question number

% -----------------------------------------------------------

\refstepcounter{qn}
\noindent
\textbf{\arabic{qn}. Si enunci il Principio di Relatività di Galileo. Siano $\vec r(t)$, $\vec v(t)$ e $\vec a(t)$ rispettivamente la posizione, la velocità e  l’accelerazione di un punto materiale in un sistema di riferimento inerziale $O$. Si esprimano posizione, velocità e accelerazione in un secondo sistema di riferimento $O’$ in moto traslatorio uniforme con velocità di trascinamento $\vec V_T$ rispetto al primo.}
\\
\\
Il Principio di Relatività di Galileo afferma che le leggi della meccanica sono le stesse in tutti i sistemi di riferimento inerziali.
\\
\\
Siano $\vec r'(t)$, $\vec v'(t)$, e $\vec a'(t)$ rispettivamente la posizione, la velocità e l'accelerazione del punto materiale rispetto a $O'$, allora
\[
	\vec {r'}(t) = \vec r(t) - \vec V_T t
\]
\[
	\vec {v'}(t) = \frac{d \vec {r'}(t)}{dt} = \frac{d}{dt}\left( \vec r(t) - \vec V_T t\right) = \vec v(t) + \vec V_T
\]
\[
	\vec {a'}(t) = \frac{d \vec{v'}(t)}{dt} = \frac{d}{dt}\left( \vec v(t) - \vec V_T \right) = \vec a(t)
\]
\vspace{1cm}

% -----------------------------------------------------------

\refstepcounter{qn}
\noindent
\textbf{\arabic{qn}. Si descriva la legge di trasformazione classica delle velocità tra sistemi di riferimento e si faccia un esempio
	pratico.}
\\
\\
Se un sistema di riferimento inerziale $O'$ si muove rispetto ad un altro sistema di riferimento inerziale $O$ con velocità di trascinamento $\vec V_T$, allora le velocità sono correlate da
\[
	\vec{v} = \vec {v'} + \vec V_T
\]
dove $\vec v$ è la velocità del corpo rispetto al sistema $O$, e $\vec{v'}$ è la velocità del corpo rispetto al sistema $O'$.
\\
\\
Supponiamo che un treno (sistema $O'$) si muova con velocità costante $\vec V_T = 50$ km/h verso est rispetto al suolo (sistema $O$). Un passeggero all'interno del treno cammina verso la parte anteriore del treno con velocità $v' = 5$ km/h rispetto al treno. Utilizzando la formula , otteniamo che il passeggero si muove a velocità 5 km/h + 50 km/h = 55 km/h rispetto al suolo nella direzione est.



% -----------------------------------------------------------
% End Document
% -----------------------------------------------------------
\end{document}

